\begin{textsample}{POS Dim 3 – pe – Score 10.00 – pe/pe\_inf\_17.txt}  \label{ex:f3_pos_190}
2.4 FUNÇÃO SOCIOPOLÍTICA E PEDAGÓGICA DA EDUCAÇÃO \textbf{INFANTIL} De acordo com o marco legal que norteia a Educação \textbf{Infantil}, a função sociopolítica dessa primeira etapa da Educação básica determina que o espaço de educação coletiva fora do contexto familiar deve se inscrever no projeto de sociedade democrática com responsabilidades no desempenho de um papel ativo na construção de uma sociedade livre, justa e solidária.
Que tenha como propósito a redução das desigualdades sociais e regionais e a promoção do bem de todos ( art. 3º, incisos II e IV da Constituição Federal ), compromissos a serem perseguidos pelos sistemas de ensino e por todos os educadores na Educação \textbf{Infantil}.
No que tange à função pedagógica da Educação \textbf{Infantil}, é parte integrante da Educação Básica, como diz a Lei nº 9.394/96 em seu artigo 22, e tem por finalidade desenvolver o educando, assegurando -lhe a formação comum indispensável para o exercício da cidadania e fornecer -lhe meios para progredir no trabalho.
Essa finalidade deve ser adequadamente interpretada em relação às \textbf{crianças} \textbf{pequenas}, compreendendo, nessa interpretação, as formas como as \textbf{crianças}, nessa etapa de vida, interagem, aprendem e se desenvolvem de forma bem específica e indicando uma identidade própria para esse segmento que, através do cuidar e o educar, rompe com a visão assistencialista e define esta concepção como direito das \textbf{crianças}.
O cuidar e educar são funções indissociáveis na Educação \textbf{Infantil}, materializados quando a instituição garante a integração dos aspectos físicos, emocionais, cognitivos e sociais nas propostas desenvolvidas, rompendo com a cultura da divisão entre os \textbf{cuidados} com a higienização e alimentação e as práticas pedagógicas.
Nesse sentido, para Barbosa ( 2006 ), o cuidar ultrapassa processos ligados à proteção e ao atendimento das necessidades físicas de alimentação, repouso, \textbf{higiene}, \textbf{conforto} e prevenção da dor.
Exige que o professor se coloque disponível para a \textbf{escuta} das necessidades, dos \textbf{desejos} e inquietações das \textbf{crianças}, no sentido de \textbf{apoiar} as suas conquistas.
As práticas envolvidas nos atos de alimentar -se, tomar banho, trocar \textbf{fraldas} e controlar os esfíncteres, na escolha do que vestir, na atenção aos riscos de adoecimento mais fácil nessa faixa etária, no âmbito da Educação \textbf{Infantil}, não são apenas práticas que respeitam o direito da \textbf{criança} de ser bem atendida nesses aspectos, como cumprimento do respeito à sua dignidade como pessoa humana.
Elas são também práticas que respeitam e atendem ao direito da \textbf{criança} de apropriar -se, por meio de experiências corporais, dos modos estabelecidos culturalmente de alimentação e promoção de saúde, de relação com o próprio corpo e consigo mesma, mediada pelos professores que intencionalmente planejam e cuidam da organização dessas práticas.

% matched lemmas: apoiar, conforto, criança, cuidado, desejo, escuta, fralda, higiene, infantil, pequeno
\end{textsample}
